\documentclass[12pt]{amsart}

\def\CHT{05} \def\SEC{01} \def\SEM{S} \def\YEAR{2018} \def\CLASS{MA311}

\newcommand{\BACKROOT}{../../../..}
\newcommand{\TEACHINGROOT}{\BACKROOT}
%\newcommand{\TEACHINGROOT}{/Users/rhcocker/Dropbox/Teaching}
\newcommand{\COURSEROOT}{\TEACHINGROOT/\SEM_\YEAR/\CLASS}
\newcommand{\SECTIONROOT}{\COURSEROOT/Chapter_\CHT/Section_\CHT_\SEC}
\newcommand{\PATHTOHEADERS}{\TEACHINGROOT/Headers}
\newcommand{\PATHTOSLIDES}{\SECTIONROOT/Slides}
\newcommand{\PATHTOFIGS}{\COURSEROOT/Figures}

\usepackage{mathrsfs}
\newcommand{\LAP}[1]{\mathscr{L}\{#1\}}

\input{\PATHTOHEADERS/TestQuizPackageHeader.tex}

\renewcommand{\headrulewidth}{0pt} % remove lines as well

\numberwithin{equation}{section}

\DeclareGraphicsRule{.tif}{png}{.png}{`convert #1 `dirname #1`/`basename #1 .tif`.png}

\textheight = 10.8 in
\topmargin = -1.0in
\textwidth = 7.6in
%\oddsidemargin = -0.7in
%\evensidemargin = -0.7in

\newcounter{ProbNum}
\setcounter{ProbNum}{0}
\def\NEWPROB{\stepcounter{ProbNum} {\noindent \arabic{ProbNum}. \,\,}}

\begin{document}

\input{\PATHTOHEADERS/TestDefinitionHeader.tex}

\pagestyle{empty}

\noindent\rule{7.6in}{0.4pt}

{\centering\sc 
\noindent Differential Equations\\
 Laplace Transforms \\
}
\noindent\rule{7.6in}{0.4pt}

\vspace{0.1cm}

\UL{Laplace Transform}

\noindent  The Laplace Transform of function $f(t)$ is a new function $F(s)$ defined by:
	\[ F(s) = \LAP{f(t)} = \int_0^{\infty}e^{-st}f(t) dt, \]
	provided this integral is finite (converges).\\
	
\NEWPROB Use the definition of Laplace Transform (above) to find the Laplace Transform of $h(t) = 7t.$  
	Include appropriate restrictions on $s$ to guarantee the convergence of the integral.\\

\newpage
\thispagestyle{empty}

\NEWPROB Use the definition of Laplace Transform (above) to find the Laplace Transform of  the function
	\[d(t) = \left\{ \begin{array}{ll}	0,	& t<0\\
							7,	& 0 \le t < 3\\
							0,	& t \ge 3\\
							\end{array} \right.\]  %= 7U(t) - 7U(t-3)
	Include appropriate restrictions on $s$ to guarantee the convergence of the integral.\\
	
%\item    $\D H(s) = \frac{7}{s^2},$ $s>0$\\
%
%\item    $\D W(s) = \frac{1}{s-k},$  $s > k$	\\
%
%\item   $\D Y(s) = \frac{2}{s^2} + \frac{1}{s-5},$  $s>5$

%\item  $\D D(s) = \frac{7}{s} - \frac{7e^{-3s}}{s}$

\newpage
\thispagestyle{empty}

\NEWPROB Show that the Laplace Transform has the following 2 properties:

\begin{enumerate}
	\item  $\D \LAP{f(t)+g(t)} = \LAP{f(t)}+\LAP{g(t)}$\\ \vspace{10.2cm}
	\item  $\D \LAP{A f(t)} = A \LAP{f(t)}\quad $($A$ is constant)\\
\end{enumerate}

\newpage
\thispagestyle{empty}

\NEWPROB Determine the Laplace transform of each function. Use the table given in class.
%-------NUM LIST---------
\begin{itemize}
	\item  $\D f(t) = e^{-2t}$\\ \vspace{5.2cm}
	\item  $\D g(t) = t^3 $\\ \vspace{5.2cm}
	\item  $\D h(t) = e^{2t} - \sin (5t)$\\ \vspace{5.2cm}
	\item  $\D q(t) = 8t + 4 - \cos(6t) $\\ \vspace{5.2cm}
\end{itemize}
		
\newpage
\thispagestyle{empty}

\end{document}